% from ega3-1-fr.tex
\oldpage[III]{154}

Comme $h$ est propre (II, 5.4.2), donc fermé, il résulte de
(\hyperref[III.5.1.3.1-fr]{5.1.3.1}) que $X\cap T=T$, autrement dit $T\subset X$,
et comme $T$ est fermé dans $Z$, $T$ est propre sur $Y$. Si $\mathcal F$ est le
faisceau induit sur $X$ par $\mathcal G$, son complété $\widehat{\mathcal F}$ est
induit sur $\widehat X$ par $\widehat{\mathcal G}$ (I, 10.8.4), donc est égal à
$\mathcal F'$, ce qui achève la démonstration.
\end{env}

\subsection{Démonstration du théorème d'existence : cas général}
\label{subsection:III.5.3-fr}

\begin{lemma}[5.3.1]
\label{III.5.3.1-fr}
Sous les conditions de (\hyperref[III.5.1.4-fr]{5.1.4}), si
\[
0\to\mathcal H\to\mathcal F\to\mathcal G\to0
\]
est une suite exacte de $\mathcal O_{\widehat X}$-Modules cohérents telle que
$\mathcal G$ et $\mathcal H$ soient algébrisables, alors $\mathcal F$ est
algébrisable.
\end{lemma}

\begin{proof}
En effet, supposons que $\mathcal G=\widehat{\mathcal B}$,
$\mathcal H=\widehat{\mathcal C}$, $\mathcal B$ et $\mathcal C$ étant des
$\mathcal O_X$-Modules cohérents à supports propres sur $Y$~; $\mathcal F$ définit
canoniquement un élément du $A$-module
\[
\operatorname{Ext}^1_{\mathcal O_{\widehat X}}
(\widehat X;\widehat{\mathcal B},\widehat{\mathcal C})
\]
($0$, 12.3.2), et les hypothèses entraînent que ce $A$-module est canoniquement
isomorphe à
\[
\operatorname{Ext}^1_{\mathcal O_X}(X;\mathcal B,\mathcal C)
\]
(4.5.2)~; il existe donc une suite exacte
$0\to\mathcal C\to\mathcal A\to\mathcal B\to0$ de $\mathcal O_X$-Modules
cohérents telle que l'image canonique de l'élément de
$\operatorname{Ext}^1_{\mathcal O_X}(X;\mathcal B,\mathcal C)$ correspondant à
$\mathcal A$ soit l'élément de
$\operatorname{Ext}^1_{\mathcal O_{\widehat X}}
(\widehat X;\widehat{\mathcal B},\widehat{\mathcal C})$ correspondant à
$\mathcal F$. Mais par définition (compte tenu de (I, 10.8.8, (iii))), cela signifie
que $\mathcal F$ est isomorphe à $\widehat{\mathcal A}$, d'où le lemme, car
$\operatorname{Supp}(\mathcal A)$ est contenu dans la réunion de
$\operatorname{Supp}(\mathcal B)$ et $\operatorname{Supp}(\mathcal C)$, donc est
propre sur $Y$.
\end{proof}

\begin{corollary}[5.3.2]
\label{III.5.3.2-fr}
Sous les conditions de (\hyperref[III.5.1.1-fr]{5.1.1}), soit
$u:\mathcal F\to\mathcal G$ un homomorphisme de
$\mathcal O_{\widehat X}$-Modules cohérents~; si $\mathcal G$,
$\operatorname{Ker}(u)$ et $\operatorname{Coker}(u)$ sont algébrisables, il en est
de même de $\mathcal F$.
\end{corollary}

\begin{proof}
Le lemme (\hyperref[III.5.2.2-fr]{5.2.2}) appliqué à l'homomorphisme
$\mathcal G\to\operatorname{Coker}(u)$ montre en effet que
$\operatorname{Im}(u)$ est algébrisable, et il suffit ensuite d'appliquer le lemme
(\hyperref[III.5.3.1-fr]{5.3.1}) à la suite exacte
\[
0\to\operatorname{Ker}(u)\to\mathcal F\to\operatorname{Im}(u)\to0.
\]
\end{proof}

\begin{lemma}[5.3.3]
\label{III.5.3.3-fr}
Sous les conditions de (\hyperref[III.5.1.1-fr]{5.1.1}), soient $h:Z\to Y$ un
morphisme de type fini, $\widehat Z$ le complété de $Z$ le long de
$Z'=h^{-1}(Y')$, $g:Z\to X$ un $Y$-morphisme propre,
$\widehat g:\widehat Z\to\widehat X$ son prolongement aux complétés. Pour tout
$\mathcal O_{\widehat Z}$-Module algébrisable $\mathcal F'$,
$\widehat g_*(\mathcal F')$ est un $\mathcal O_{\widehat X}$-Module algébrisable.
\end{lemma}

\begin{proof}
En effet, si $\mathcal F$ est un $\mathcal O_Z$-Module cohérent tel que
$\mathcal F'=\widehat{\mathcal F}$, il résulte du premier th. de comparaison
(4.1.5) que $\widehat g_*(\widehat{\mathcal F})$ est isomorphe au complété de
$g_*(\mathcal F)$.
\end{proof}

\begin{lemma}[5.3.4]
\label{III.5.3.4-fr}
Soient $X$ un schéma (usuel) noethérien, $X'$ une partie fermée de $X$,
$f:Z\to X$ un morphisme propre, $Z'=f^{-1}(X')$,
$\widehat X=X_{/X'}$, $\widehat Z=Z_{/Z'}$,
$\widehat f:\widehat Z\to\widehat X$ le prolongement de $f$ aux complétés. Soit
$\mathcal M$ un Idéal cohérent de $\mathcal O_X$ tel que, si
$U=X-\operatorname{Supp}(\mathcal O_X/\mathcal M)$, la restriction
$f^{-1}(U)\to U$ de $f$ soit un isomorphisme. Alors, pour tout
$\mathcal O_{\widehat X}$-Module cohérent $\mathcal F$, il existe un entier $n>0$
tel que le noyau et le conoyau de l'homomorphisme canonique
\[
\rho_{\mathcal F}:\mathcal F\longrightarrow
\widehat f_*\bigl(\widehat f^{,*}(\mathcal F)\bigr)
\]
($0_{\mathrm I}$, 4.4.3) soient annulés par $\widehat{\mathcal M}^{,n}$.
\end{lemma}

\begin{proof}
On peut se borner au cas où $X=\operatorname{Spec}(B)$ où $B$ est un anneau
noethérien, donc $X'=\mathrm V(\mathfrak K)$, où $\mathfrak K$ est un idéal de
$B$. Nous allons voir qu'on peut se ramener au cas où $B$ est un anneau adique
noethérien et $\mathfrak K$ un idéal de définition de $B$. Soit en effet $B_1$ le
séparé complété de $B$ pour la topologie $\mathfrak K$-préadique~; si
$\mathfrak K_1=\mathfrak K B_1$, $B_1$ est donc un

